\documentclass[12pt,a4paper]{report}
\usepackage{amsmath}
\usepackage{amssymb}
\usepackage[backend=biber]{biblatex}
\addbibresource{Project.bib}
\usepackage{parskip}
\setlength{\emergencystretch}{3em}
\begin{document}
\chapter{Project Proposal}
\begin{tabular}{@{}l@{ : }l}
Name & \bf Oluwatosin Bernice (Tosin Bea) Jokosenumi\\
College & \bf Churchill College \\ 
CRSID & \bf obaj2 \\
Project Supervisors & \bf Chiara D'Ercoli, Amogh Singh \\
Director of Studies & \bf Dr John Fawcett \\ 
Project Checkers & \bf Professor Pietro Lio' \\
\end{tabular}
\section{Introduction and Description}

Multiple sclerosis is a polygenic, complex autoimmune disorder where the immune system attacks myelin, the protective cover of nerve cells (axons) in the central nervous system (CNS). Although damage can  be  mitigated to an extent by the re-routing of messages using undamaged nerve cells (plasticity), the damage eventually becomes too large to compensate resulting in progressive neurodegeneration.

Gene expression is the process outlined by the central dogma of molecular biology: DNA is transcribed into RNA, and RNA is translated into functional products. The measure of how much a gene is expressed at a given time is referred to as its gene expression level, and the analysis of differential expression levels when comparing healthy people and diagnosed patients can aid in the development of disease signatures based on changes in gene activity. Advanced, high-throughput technologies for studying gene expression, such as single-cell RNA sequencing (scRNA-seq) and single-nucleus RNA sequencing (snRNA-seq), allow gene expression to be measured at the level of individual cells or nuclei, enabling the heterogeneity of cell populations to be explored.  Studies have utilised snRNA-seq data in clustering distinct  MS sub-groups \parencite{macnairSingleNucleiRNAseq2022} providing a basis for use in future individualised treatment and prediction.

Given that the complexity of biological networks in higher organisms  primarily arises from dynamic gene interactions as opposed to the size of the genome (C-value paradox), analysing gene dynamics offers a more informative framework for understanding biological systems. Gene-gene co-expression networks (GGCNs) are typically presented as undirected graphs where nodes are genes and the presence of edges depends on co-expression values calculated from gene expression data.\parencite{vandamGeneCoexpressionAnalysis2017} These networks can be used to  identify the functional status of genes from a systematic perspective, which can in turn inform the discovery of gene-disease associations through guilt-by-association analyses.

Graph neural networks (GNNs) have commonly been used for tasks like graph classification, node classification and link prediction.\parencite{sewellGNNsNetworkClassification2025} Given a graph $G(V,E)$  as input in the form of a node feature matrix ($H:  \mathbb{R}^{N \times F}$),  GNNs produce new node embeddings that capture the node features, and the learnt embeddings can be used for graph classification. Graph Attention Networks (GATs) introduce a masked attention mechanism which take the neighbourhood of a node into account in computing the updated features for a node\parencite{velickovicGraphAttentionNetworks2017} - previous works have successfully applied these networks to disease prediction\parencite{ravindraDiseaseStatePrediction2020}, providing precedent for its use in the project.

The project is composed of two main goals:
\begin{itemize}
    \item Using the hdWGCNA framework to construct multi-layer co-expression networks from snRNA-seq gene-expression data

\item Constructing a Multiplex Graph Attention Network (mGAT) which takes the GGCNs as input and outputs a binary classification for MS with a structure-blind multilayer perceptron (MLP) model and a vanilla Graph Convolution Network (GCN) to serve as baseline comparisons\parencite{luzhnicaGraphClassificationNetworks2019}.
\end{itemize}
\section{Starting Point}
I have some experience with machine learning beyond the Computer Science Tripos, however I was  unfamiliar with graph attention networks, so I have done some preliminary reading. In particular, I have read [1-6] however further reading will be required as noted in the timetable. 

\section{Success Criterion and Evaluation}
The project will be considered a success if the following criteria are met:

\begin{itemize}
    \item A structure blind MLP, mGAT and vanilla GCN are designed and implemented for graph classification 
    \item Real-world transcriptomic data is used to train the model
    \item The evaluation framework is carried out on the models to assess their performance
\end{itemize}

The evaluation framework for the project will use a confusion matrix, allowing for 4 classification metrics to be calculated: accuracy, precision, recall and a recall-oriented F-score given that missing a diagnosis of MS may delay medical intervention. These metrics are defined as, given the true positives (TP), true negatives (TN), false positives (FP) and false negatives (FN) from the confusion matrix, as:
\begin{align*}
    Precision &= \frac{TP}{TP+FP} \\
    Recall &= \frac{TP}{TP+FN} \\
    F_\beta-score &=(1+ \beta^2) \frac{Precision \times Recall}{(\beta^2 \times Precision)+Recall} \\
    Accuracy &= \frac{TP+TN}{TP+TN+FP+FN}
\end{align*}

\section{Extensions and Further Work}
If the success criteria are met, I will attempt to implement the following extensions: 
\begin{itemize}
    \item \textbf{EX1}: Incorporate epigenomic (snATAC-seq) data to produce a multi-omic classification model. \\
\item\textbf{EX2}: Implement a multi-class classification model using transcriptomic and epigenomic data from Alzheimers' patients and evaluate with a new evaluation framework.\\
\item\textbf{EX3}: Improve the explainability of the model by evaluating feature importance.
\end{itemize}
\section{Work Plan}
\textbf{Michaelmas Term} \\ 
\begin{itemize}
    \item
    \textbf{Weeks 1-2 (13/10-26/10)} \\
    \begin{itemize}
        \item Background reading on mGATs, GCNs and the hdWGCNA package and set up code repository.
        \item Milestone: Produce a document explaining the theory about these model and the reason for their inclusion in the project.
    \end{itemize}
    \item \textbf{Weeks 3-4 (27/10-9/11)} \\
    \begin{itemize}
        \item Read up on PyTorch and decide on the network architecture for the three models.
        \item Milestone: Have diagrams for each network architecture and produce basic ML examples in PyTorch.
    \end{itemize}
    \item \textbf{Weeks 5-6 (10/11-23/11)} \\
    \begin{itemize}
        \item Process the snRNA-seq dataset using the py-hdWGCNA framework and split data into training, validation and testing datasets at a patient level
        \item Milestone: The datasets should be ready for training 
    \end{itemize}
    \item \textbf{Weeks 7-8 (24/11-7/12)} \\
    \begin{itemize}
        \item Buffer period
    \end{itemize}
\end{itemize}
\textbf{Winter Break} \\ 
\begin{itemize}
    \item \textbf{Weeks 9-10 (8/12-21/12)} \\
    \begin{itemize}
        \item 
    \end{itemize}
    \item \textbf{Weeks 11-12 (22/12-4/1)} \\
    \item \textbf{Weeks 13-14 (5/1-18/1)} \\
\end{itemize}
\textbf{Lent Term}
\begin{itemize}
    \item \textbf{Weeks 15-16 (19/1-1/2)} \\
    \item \textbf{Weeks 17-18 (2/2-15/2)} \\
    Friday 5 February 12pm: Progress Report Presentation
    \item \textbf{Weeks 19-20 (16/2-1/3)} \\
    Friday 26 February: Progress Report Presentations
    \item \textbf{Weeks 21-22 (2/3-15/3)} \\
\end{itemize}
\textbf{Easter Break} \\ 
\begin{itemize}
    \item \textbf{Weeks 23-24 (16/3-29/3)} \\
    \item \textbf{Weeks 25-26 (30/3-12/4)} \\
    \item \textbf{Weeks 27-28 (13/4-26/4)} \\
\end{itemize}

\textbf{Lent Term}
\begin{itemize}
    \item \textbf{Weeks 15-16 (27/4-9/5)} \\
    Friday 11 June 12pm: Dissertation Deadline 
\end{itemize}
\section{Resource Declaration}

I will be using my personal laptop, a 2024 Dell Inspiron 16 2-in-1, with 16GB memory and a 1TB SSD for software development. As a backup, I will use my secondary personal laptop with an external 1TB storage. I will continuously backup my code and dissertation with Git version control. I accept full responsibility for this machine and I have made contingency plans to protect myself against hardware and/or software failure. In terms of software, I will use Python including the PyTorch and py-hdWGCNA libraries. 

The project also requires the use of biological datasets for training data. These datasets have been taken from https://zenodo.org/records which are publicly available to anyone.

\printbibliography
\end{document}